%% -*- coding: utf-8 -*-
%% EXAMPLE CONTENT — ships with a new SciTeX workspace. Replace it with yours.
\section{Introduction}
\pdfbookmark[1]{Introduction}{introduction}

\textbf{This is an example, not a study.} It is shipped with every new SciTeX
workspace so that the editor, the figure browser, and the file tree have real
content in them the first time you open them. Nothing below is a novel finding,
and the manuscript should not be cited.

An empty template is a poor way to learn a tool. It shows you where things go
but not what they are for, and it gives you nothing to change. The alternative
is a project small enough to read in one sitting, honest enough that every
number in it can be re-derived, and boring enough that no one mistakes it for a
claim. Handwritten digit recognition is the standard choice for exactly those
reasons: the task is easy to state, the data is public, and the expected result
has been known for three decades.

The specific question is narrow. Optical character recognition normally begins
by extracting features --- strokes, moments, edges --- before any classifier
sees the image. What happens if that step is skipped entirely, and a linear
decision rule is applied directly to the raw pixel intensities of a very small
thumbnail? The images used here are $8\times8$, coarse enough that a human
reader can still name most digits but fine detail is gone.

The answer is that most of the work is already done. A linear support vector
machine on 64 raw pixel values recovers the great majority of the labels, and
its remaining errors are concentrated in the one place a linear rule is least
able to help: the digits whose ink occupies nearly the same pixels. Section
\ref{sec:methods} describes the data and the fit, Section \ref{sec:results}
reports what it does, and Section \ref{sec:discussion} says what that does and
does not show.

More usefully for a new user, this project also shows the shape of a SciTeX
workspace. The data lives in \texttt{data/}, one seeded script in
\texttt{scripts/} turns it into the figures in \texttt{figures/}, and this
manuscript reads those figures back. Change the CSV and re-run the script, and
the numbers below change with it.

%%%% EOF
